\section{Derivations}
\label{app:deriv}

\subsection{The singular box}

The in-support integral of Eq.~\eqref{eq:Fc} is split by writing the kernel
as its divergence plus a smooth remainder. The divergence part is the
eight-term $\mathrm{asinh}$ box
\begin{equation}
\begin{split}
T(A,B,u)&=\int_0^A\!\int_0^B\frac{\mathrm{d}X\,\mathrm{d}Y}
{\sqrt{u^2+X^2+Y^2}}\\
&=\begin{cases}
A\,\mathrm{asinh}(B/A)+B\,\mathrm{asinh}(A/B), & |u|<10^{-300},\\[2pt]
A\,\mathrm{asinh}\!\bigl(B/\sqrt{u^2+A^2}\bigr)
+B\,\mathrm{asinh}\!\bigl(A/\sqrt{u^2+B^2}\bigr)\\
\qquad
-u\,\arctan\!\bigl(AB/(u\sqrt{A^2+u^2+B^2})\bigr), & \text{otherwise.}
\end{cases}
\end{split}
\end{equation}
with $A,B>0$ (Eq.~\eqref{eq:Tprimitive}). Each box face has $8$ corners and
$24$ such calls per component, so the eight-term box is $\mathrm{O}(1)$ per
pair. The form is even in $u$: $u$ is the distance from the field point to
the face plane, so a field point beyond a box face gives negative $u$, and
the $|u|\to0$ branch must fire on $|u|$, not $u$. The branch at
$|u|<10^{-300}$ is the $u\to0$ limit, which is finite; it is not a
regularization, it is the limit of the second branch as $|u|\to0$.

The primitive is verified against a converged midpoint trapezoid reference:
the double integral of the integrand over the unit box with $u=1$ is
$-2.6348$ for the converged midpoint trapezoid; the naive trapezoid gives
$-2.6333$, $1.5\times10^{-3}$ low, a trapezoid singularity artefact at the
corner where the integrand is not smooth. The primitive is implemented in
the fast evaluator header \texttt{include/picBunch/detail/FastNear.hpp} and
the device reference
\texttt{tools/closed-form/gpu\allowbreak \_extract.cu}.

\subsection{The smooth remainder}

The remainder integrates $(S-W^\ast)K$ over the TSC knot panels, where $S$
is the step function of the kernel's support and $W^\ast$ is the TSC shape
subtracted so that the integrand is smooth inside the support. Each axis is
split at the TSC knots ($-1.5, -0.5, 0.5, 1.5$) and at the anchor point
$s^\ast$, the projection of the source onto the axis; on each panel a plain
tensor Gauss--Legendre quadrature is applied. The default order is $5$ per
axis, $8000$ nodes per in-support pair. For a moving bunch the box is
$\gamma$-stretched for the axial direction: the beaming cone compresses the
field into a pancake of width $d/\gamma$, and the stretched box matches
that compression so the singular part stays $\mathrm{O}(1)$. The oblique
direction uses a rest-frame octant quadrature of the singular box at order
$8$ ($4096$ nodes) plus the same remainder, because the octant structure of
the rest-frame kernel does not map to the stretched box.

\subsection{The $h^3$ Riemann factor}

The sampled convolution evaluates the continuum integral of
Eq.~\eqref{eq:Fc} by a Riemann sum whose nodes are spaced $h$ cells apart,
so each sample covers a physical volume $h^3$ (in cell units). The
unnormalized FFT round trip scales by $M^3$, so the read-out applies the
factor $h_x h_y h_z/(M_x M_y M_z)$ to convert the sum to the integral.
This is the only place where the mesh spacing enters the amplitude; the
scale-invariance of Section~\ref{sec:scaleinvariance} follows from it,
together with the homogeneity of degree $-2$ of the kernel and the cell-unit
Yee offsets.
